% ============================================
% METHODS - Showcase: tables, blocks, enumerations, math
% ============================================

% --------------------------------------------
% Slide: Comparison Table with Checkmarks
% --------------------------------------------
\begin{frame}{Comparison of Approaches}
\framesubtitle{Why existing methods fall short}

\begin{table}
\centering
\small
\begin{tabular}{lccc}
\toprule
\textbf{Property} & \textbf{Method A} & \textbf{Method B} & \textbf{Ours} \\
\midrule
Scalable & \cmark & \xmark & \cmark \\
Interpretable & \xmark & \cmark & \cmark \\
Low latency & \cmark & \cmark & \cmark \\
Domain-agnostic & \xmark & \xmark & \cmark \\
\bottomrule
\end{tabular}
\end{table}

\vspace{0.5em}

\begin{colorblock}[white]{stevensred}{Advantage}
Our approach is the first to satisfy all four desirable properties simultaneously.
\end{colorblock}

\end{frame}

% --------------------------------------------
% Slide: Pipeline / Enumeration
% --------------------------------------------
\begin{frame}{Proposed Pipeline}
\framesubtitle{Three-stage processing}

\begin{columns}
\begin{column}{0.48\textwidth}
\textbf{Stage 1: Data Preprocessing}
\begin{enumerate}
    \item Collect raw data from source
    \item Clean and normalize
    \item Split into train/val/test
\end{enumerate}

\vspace{0.5em}
\textbf{Stage 2: Model Training}
\begin{enumerate}
    \item Fine-tune pre-trained model
    \item Hyperparameter search
    \item Early stopping on validation set
\end{enumerate}
\end{column}
\begin{column}{0.48\textwidth}
\textbf{Stage 3: Evaluation}
\begin{enumerate}
    \item Run on held-out test set
    \item Compute metrics (P, R, F1)
    \item Statistical significance tests
\end{enumerate}

\vspace{0.5em}

\begin{colorblock}[white]{stevensred}{Key Insight}
Pre-training captures domain knowledge; fine-tuning adapts with minimal supervision.
\end{colorblock}

\end{column}
\end{columns}
\end{frame}

% --------------------------------------------
% Slide: Block Types (standard, alert, example)
% --------------------------------------------
\begin{frame}{Block Styles}
\framesubtitle{Three built-in block types}

\begin{block}{Standard Block}
This is a standard block with the Stevens red header. Use it for general information, definitions, or neutral content.
\end{block}

\vspace{0.3em}

\begin{alertblock}{Alert Block}
This is an alert block. Use it for warnings, limitations, or important caveats that readers should be aware of.
\end{alertblock}

\vspace{0.3em}

\begin{exampleblock}{Example Block}
This is an example block with a green header. Use it for examples, case studies, or positive results.
\end{exampleblock}
\end{frame}

% --------------------------------------------
% Slide: Math and Formulas
% --------------------------------------------
\begin{frame}{Mathematical Formulation}

\textbf{Objective function:}
\[
\mathcal{L}(\theta) = -\frac{1}{N}\sum_{i=1}^{N} \log p_\theta(y_i \mid x_i) + \lambda \|\theta\|_2^2
\]

\vspace{0.3em}
\textbf{Where:}
\begin{itemize}
    \item $\theta$ --- model parameters
    \item $(x_i, y_i)$ --- input-output pairs
    \item $\lambda$ --- regularization strength
\end{itemize}

\vspace{0.5em}

\begin{colorblock}[white]{stevensred}{Note}
The regularization term prevents overfitting when training data is limited.
\end{colorblock}

\end{frame}
